% vim: set ft=tex tabstop=4 shiftwidth=4 noexpandtab:

\input{../../include/latex/tikz/header.tex}

% opening %{{{1

\begin{document}
\begin{tikzpicture}[scale=1.0]

% parameters %{{{1

	\tikzmath{
		\radius = 3.2 ;
		\rotation = 90 ;
		\numsidesfloat = 5.0 ;
		\numsides = int(\numsidesfloat) ;
		int \i ;
		for \i in {1,...,\numsides} {
			\ang{\i} = \rotation + 360 / \numsidesfloat * (\i - 1) ;
		};
	}

% coordinates %{{{1

	\coordinate (O) at (0,0);

	\foreach \i in {1,...,\numsides}
		\coordinate (P_\i) at (\ang{\i}:\radius);

% points, dots, vertices %{{{1

	\fill (O) circle (0.4mm);

	\foreach \i in {1,...,\numsides}
		\fill (P_\i) circle (0.4mm);

% segments, sides, lines %{{{1

	\draw (P_1) \foreach \i in {2,...,\numsides} { -- (P_\i) } -- cycle;

	\draw (P_1) -- (P_3);
	\draw (P_1) -- (P_4);

% circles %{{{1

	\draw (O) circle (\radius);

	%\path[draw] let \p1 = ($ (A) - (O) $)
	%in (O) circle ({veclen(\x1,\y1)});

% points, dots, vertices labels %{{{1

	\node[anchor=north] at ($ (O) + (220:0.0) $) {$O$};

	\node[above] at (P_1) {$P_1$};
	\node[left] at (P_2) {$P_2$};
	\node[below left] at (P_3) {$P_3$};
	\node[below right] at (P_4) {$P_4$};
	\node[right] at (P_5) {$P_5$};

% segments, sides, lines labels %{{{1

	\node[above left] at ($ (P_1)!0.5!(P_3) $) {$d_1$};
	\node[above right] at ($ (P_1)!0.5!(P_4) $) {$d_2$};

% segments, sides, lines marks %{{{1

		\foreach \i
			[evaluate=\i as \j using { int(mod(\i,\numsides) + 1) }]
			in {1,...,\numsides}
		{
			\tkzMarkSegments[mark=||, size=2pt](P_\i,P_\j);
		}

% circles labels %{{{1

	\node at ($ (O) + (60:\radius + 0.3) $) {$\mathscr{C}$};

% angles labels %{{{1

	\pic[draw, ->, "$108^\circ$", angle radius=0.5cm, angle eccentricity=1.8]
		{angle = P_3--P_2--P_1};

	\pic[draw, ->, "$36^\circ$", angle radius=1.3cm, angle eccentricity=1.3]
		{angle = P_2--P_1--P_3};

	\pic[draw, ->, "$36^\circ$", angle radius=1.3cm, angle eccentricity=1.3]
		{angle = P_1--P_3--P_2};

	\pic[draw, ->, "$108^\circ$", angle radius=0.5cm, angle eccentricity=1.8]
		{angle = P_1--P_5--P_4};

	\pic[draw, ->, "$36^\circ$", angle radius=1.3cm, angle eccentricity=1.3]
		{angle = P_4--P_1--P_5};

	\pic[draw, ->, "$36^\circ$", angle radius=1.3cm, angle eccentricity=1.3]
		{angle = P_5--P_4--P_1};

% closing %{{{1

\end{tikzpicture}
\end{document}
